\documentclass[10pt,twocolumn]{article}
\usepackage[a4paper,margin=0.7in]{geometry}
\usepackage[T1]{fontenc}
\usepackage{lmodern,amsmath,booktabs,graphicx,hyperref}
\title{When Do Utility-Aware ANN Search Budgets Really Save Work?\\
\large The Role of Relevance-Judgment Completeness and Risk Calibration}
\author{Hrithik Sharma\\Indian Institute of Technology Delhi}
\date{Research draft --- 16 November 2026}
\begin{document}
\maketitle
\begin{abstract}
We reexamine whether downstream-utility-aware ANN budget allocation provides an efficiency advantage after accounting for judgment sparsity and fixed-budget calibration.
Saved million-scale HNSW runs show that exact-positive queries contribute 92.59\% of CURA's work reduction versus the original calibration-selected fixed baseline, ruling out a purely zero-reference explanation.
However, a post-hoc fixed efSearch of 48 matches CURA's observed 2/200 utility violations and uses 28.29\% less distance work. On 43 TREC-DL-2019 queries a fixed setting also outperforms the adaptive policy at its observed risk, while ANN-safe and IR-safe oracle budgets scarcely differ.
These results support an analysis of metric-dependent headroom and calibration rather than a claim that CURA generally beats fixed HNSW.
Recovering original TREC ranked IDs enabled a same-ranking experiment with 50 prepared thinning seeds. Reducing available positive judgments increases measured ANN-minus-IR safe-budget headroom from $-0.140$ under FULL to $1.956$ with one positive per query. Exact-reference nDCG becomes zero in 82.88\% of one-positive query--seed cases, making the reference-loss criterion vacuous there. This is sensitivity to missing judgments, not improved retrieval or actual relevance preservation.
\end{abstract}

\section{Question and scope}
When do downstream-utility-aware ANN search budgets provide genuine efficiency gains, and how are those gains affected by relevance-judgment completeness and risk calibration?
We distinguish three issues: whether ANN and IR safety require different budgets; whether a learned adaptive implementation realizes that opportunity efficiently; and whether measured opportunity survives more complete judgments.
The evaluated CURA implementation selects a budget after an efSearch=10 probe and runs a fresh search if needed. It is two-stage budget allocation, not integrated graph-traversal early exit.

\section{Saved-run protocol}
The controlled corpus contains 1,000,629 passages. The original 1,000 MS MARCO queries have query-disjoint training/calibration/test sets of 600/200/200. All 43 TREC queries are OOD. We reuse the original runs and frozen policies without retrieval or retraining.
The tested efSearch values are 10,16,24,32,48,64,96,128,192,256.
Let $N_q^*$ and $N_q(e)$ denote exact and approximate nDCG@10. Utility loss is
\[
L_q(e)=\max(0,N_q^*-N_q(e)),
\]
and a violation is $L_q(e)>0.025$. ANN safety is exact-neighbor Recall@10 at least 0.9. Unjudged documents receive relevance zero; nDCG uses exponential gain, logarithmic discount, and condition-specific ideal ranking from the full query qrels.
At $N_q^*=0$, nonnegative nDCG implies zero loss automatically. This is a limitation of the reference-based criterion.

Distance work includes the CURA probe plus the selected search, except that a selected ef10 uses only the probe. Saved timing estimates sum the corresponding search components and omit controller inference and feature processing. Their p95 is across per-query mean estimates, not end-to-end request latency.
Historical saved-run intervals use paired-query percentile bootstrap intervals with 20,000 resamples. They condition on the saved policies and omit policy-selection uncertainty. New density intervals instead resample 50 paired seed-level effects conditional on the same 43 queries and rankings; these are not query-population intervals.

\section{Zero-reference explanation}
The held-out set has 41 exact-zero queries (20.5\%) and 159 exact-positive queries (79.5\%). Relative to calibration-selected fixed ef96, work savings are 8.56\% (95\% CI $[-15.97,30.89]$) in the zero group and 23.09\% ($[13.91,31.58]$) in the positive group. Positive queries contribute 92.59\% of total work reduction.
Thus the reduction is not concentrated in the trivial exact-zero regime. This does not rule out sparse-judgment effects on exact-positive queries.
Within the positive stratum, CURA uses 2730.667 distances/query versus 3550.352, a paired reduction of 819.686 ($[494.440,1119.157]$). Violations are 2/159 versus 1/159, and mean nDCG is 0.784104 versus 0.786491. Three discordant violation pairs are insufficient to establish equal risk. The work sign-flip test has MC $p\approx0.00005$ under sign-exchangeability; all tests are exploratory.

\section{Fixed-budget risk/work frontier}
The original fixed baseline ef96 was the first candidate whose two-sided 95\% Wilson upper endpoint was at most 5\% on calibration. At ef64, 4/200 violations yield upper endpoint 5.0288\%, narrowly failing; ef96 has 1/200 and upper endpoint 2.7774\%. Pointwise bounds do not establish selection-adjusted coverage.

\begin{table}[t]
\centering\small
\begin{tabular}{lrrr}
\toprule
Method & Violations & Mean work & nDCG\\
\midrule
Fixed ef24 & 3/200 & 1182.785 & .620581\\
Fixed ef48 & 2/200 & 1956.345 & .623479\\
Fixed ef64 & 1/200 & 2448.645 & .625260\\
Fixed ef96 & 1/200 & 3431.770 & .625260\\
CURA-Full & 2/200 & 2727.965 & .623363\\
\bottomrule
\end{tabular}
\caption{Selected points from the complete ten-budget test frontier. All fixed settings are in the accompanying CSV. ef48 is a post-hoc diagnostic; ef96 is the original calibration-selected baseline.}
\end{table}

\textbf{Negative result.} Fixed ef48 is the cheapest setting with observed test risk no greater than CURA's 1\%. It uses 28.29\% less distance work ($[20.07,35.39]$), a paired difference of 771.620 distances/query ($[488.510,1078.553]$). Fixed ef64 also has lower observed risk and work than CURA.
Therefore the reported 20.51\% CURA improvement against ef96 does not demonstrate that query adaptation is needed.
The observed-risk comparator was identified on test outcomes. It is not a prospectively calibrated replacement or a risk-equivalence result; its conditional bootstrap interval excludes uncertainty in choosing a setting from the frontier.
CURA's selected-search work alone is 2054.555, plus 673.410 extra charged probe work. A smaller mean ef is not a reliable measure of implemented efficiency.

\section{Oracle opportunity and judgment regime}
On all 819 exact-positive MS MARCO queries, mean first ANN-safe and IR-safe ef are 21.485 and 13.206, yielding gap 8.278 ($[6.222,10.527]$). IR permits lower ef for 25.89\%, with aggregate ef reduction 38.53\%.
The saved oracle substitutes ef256 when no safe candidate is found. Four ANN and three IR thresholds are absent in this positive stratum. Restricting to 812 jointly feasible pairs gives gap 7.611 ($[5.933,9.387]$), so the finding persists without these fallbacks. First-safe refers only to the tested grid; ef reduction is not measured distance reduction.

All 43 TREC queries are exact-positive under FULL qrels. Mean ANN-safe and IR-safe ef are 12.279 and 12.419, gap $-0.140$ ($[-1.163,0.884]$). IR allows a smaller budget for 2/43; 38 budgets agree and 3 require more IR effort. ANN is safe at the minimum tested ef10 on 36/43 queries; IR on 35/43. Five nonzero gaps provide limited inferential support and the grid floor restricts possible gains.
For saved TREC outcomes, fixed ef32 matches CURA's 1/43 violations and uses 1109.326 distances/query versus 2384.581, or 53.48\% less ($[40.13,63.06]$). This is also post hoc.
TREC has 95.395 positive judgments/query on average versus 1.071 on MS MARCO. Different query distributions, corpus coverage and only 43 TREC queries confound a causal judgment-density interpretation.

\section{Same-ranking judgment thinning}
The original NIST passage qrels contain 9,260 judgments for the saved 43 query IDs. An external archive supplied original exact top100 indices/scores, HNSW top10 indices at every tested ef, and the original ordered 1,000,629-document map. Seven source/config/query/run files match byte-for-byte. The mandatory FULL gate reproduces all 43 exact and 430 HNSW query--budget metrics: nDCG/loss tolerance $10^{-6}$, ANN recall tolerance $10^{-9}$, and identical violation indicators. No retrieval, embedding generation, index construction or controller training was performed. ANN score evidence was reconstructed from original cached vectors without reordering IDs. No independently recorded historical map checksum or immutable encoder revision exists; archive lineage and FULL agreement support consistency rather than independent identity proof.

FULL retains all available pooled qrels; it is not complete relevance truth. All 50 prepared paired seeds retain nested per-query prefixes of $\max(1,\lceil pn\rceil)$ for $p=.5,.25$ and one positive respectively. Grades and nonpositive judgments remain unchanged. Removed positives become unjudged zero, and IDCG is recomputed from each condition's retained qrels. Scoring produces 64,930 query--budget rows with fixed rankings, ANN recall, search work and CURA choices. All 43 queries remain jointly feasible in every condition and seed.

\begin{table}[t]
\centering\small
\begin{tabular}{lrrrr}
\toprule
Qrels & Gap & IR ef & Zero\% & Lower\%\\
\midrule
FULL & $-.140$ & 12.419 & 0.00 & 4.65\\
50\% & .308 & 11.971 & 1.67 & 6.33\\
25\% & .848 & 11.431 & 9.49 & 8.42\\
One positive & 1.956 & 10.323 & 82.88 & 13.40\\
\bottomrule
\end{tabular}
\caption{Condition means over the fixed 43 queries and, for thinned conditions, 50 seeds. Gap is ANN-safe minus IR-safe ef; ANN-safe mean remains 12.279. Zero is the exact-zero fraction; Lower is the rate at which IR permits a lower budget. ef differences are not measured work savings.}
\end{table}

\textbf{Outcome A: apparent headroom increases.} Paired gap changes relative to FULL are .447 ($[.300,.593]$), .988 ($[.820,1.153]$), and 2.096 ($[2.001,2.180]$) ef units. Mean effects increase with thinning; this does not assert monotonicity for every query or seed. The 82.88\% exact-zero fraction in the one-positive condition is central: zero exact nDCG automatically satisfies the reference-loss criterion, providing no evidence of relevance preservation by itself.

Twelve density follow-ups cover gap, exact-zero fraction, IR-safe budget and IR-allows-lower rate. Unadjusted percentile CIs and two-sided Monte Carlo sign-flip diagnostics use 20,000 seed resamples/draws; Holm adjustment covers these 12 comparisons only. Sign-flip inference assumes effects symmetric about zero under the null, an assumption not established by thinning alone. IR-budget effects are exact negatives of gap effects because ANN budgets and the feasible cohort are fixed; their inference is reused and is not independent corroboration. The 64,930 rows are not independent inferential units.

\section{Conclusion and limitations}
Oracle utility-aware headroom survives removal of exact-zero MS MARCO queries, but this CURA implementation does not improve the observed fixed-budget frontier. Conservative calibration and adaptive overhead materially affect interpretation. We retain this negative result.
Controlled thinning demonstrates sensitivity to missing judgments on these fixed queries, without identifying the cause of cross-dataset differences. Sparse judged utility is not actual relevance preservation; pointwise Wilson intervals are not universal guarantees; saved search timings are not full-service performance. The experiment uses a controlled approximately 1M MS MARCO subset, one primary encoder and HNSW only, with 43 TREC OOD queries. ANN-Risk is not official DARTH, and Patience-Inspired is not official Patience. The current contribution is an analysis of when adaptivity must be justified, not a universal adaptive retrieval win.
\paragraph{Reproducibility.} The repository's \texttt{sigir\_analysis} directory contains recovered ranked IDs, provenance, FULL reproduction, all 50-seed results, statistical tables, four density figures and validation scripts. Historical completed analyses were preserved. The separate FiQA human-review record is complete at 50/50, with category and evidence-note decisions supported by the saved human decision log; preliminary model suggestions alone were not counted as human reviews. Its final taxonomy and qualitative report cover fixed HNSW ef=10 failures, separately from the model-assisted MS MARCO/TREC taxonomy. This draft has not been checked against the official ACM format or compiled for final submission.
\end{document}
